\section{Execution configuration and statistics}
\label{sec:config}

\paragraph{Frozen configuration.}
Cluster and serving configuration are part of a cell's experimental
identity, and every row records the full pinned set: reader and writer
\texttt{gpt-oss-120b} on one cluster, low reasoning effort,
temperature 0, \texttt{max\_tokens} 512, 16 concurrent workers,
SFR-Embedding-\allowbreak Mistral at dimension 4096, and the benchmark's own tokenizer. A
change to any of them defines a new arm rather than additional data. Three
settings were fixed by measurement. An empty or truncated completion raises
and the cell is re-run, since scoring it as zero would push $\Delta U$
negative and manufacture harmful memories. Concurrency sits at a measured
saturation ceiling, with the ladder pushed past the chosen value. Responses
are cached on all twelve fields of a cell's identity, verified by varying one
field and confirming a cache miss. The embedding model's determinism was
checked directly: the cosine of identical texts is exactly 1.0.

\paragraph{Noise floors.}
\label{sec:floors}
$F_{\mathrm{exec}}$ is outcome variation across seeds with everything else
fixed, measured per benchmark with a bootstrap interval and within a single
execution session. $F_{\mathrm{empty}}$ is the variation of the no-memory
arm, available at no extra cost because that arm is identical for every
memory paired with a task. $F_{\mathrm{judge}}$ is scorer disagreement on
byte-identical outputs, from re-scoring stored generations; we store full
output text rather than a digest so that this re-scoring is possible. The
floor is defined on the $\Delta U$ scale, since
$\mathrm{var}(\Delta U) = \mathrm{var}_{\mathrm{with}} +
\mathrm{var}_{\mathrm{without}} - 2\,\mathrm{cov}$; per-arm floors are
reported alongside. One hazard would fabricate the whole result: the
\textsc{without} arm's inputs are identical across memories, so a cache
serving it would reuse one generation everywhere and let every $\Delta U$
clear a zero floor. That arm is re-generated per (task, replicate) rather than served from the
response cache; in the screened run all $\lfScrRealCalls$ calls reached the
endpoint, with $\lfScrCacheHits$ cache hits.

\paragraph{Labels and flips.}
A cell is robust useful when $\mathrm{CI}_{\mathrm{lower}} > 0$ and
$\overline{\Delta U} > F$, robust harmful when $\mathrm{CI}_{\mathrm{upper}}
< 0$ and $|\overline{\Delta U}| > F$, and neutral otherwise, so a cell can be
statistically non-zero and still neutral. A memory shows a robust sign flip
only if two of its cells independently meet the useful and harmful criteria.
Every flip rate carries its denominator (memories on at least two tasks) and
a Wilson interval beside the naive rate, and a cell needs at least two
replicates to enter one. Because replicates agree on every primary cell, a
robust label reduces in practice to the sign of $\Delta U$, and flip rates are
comparable only at matched cells per memory (Section~\ref{sec:negative}).

\paragraph{The additive null.}
A flip is a property of $\Delta U$'s sign; interaction is a property of its
variance. Under $\Delta U(m,t) = \mu + a_m + b_t$ with no interaction, a
memory still helps one task and harms another whenever $b_t$ crosses
$-(\mu + a_m)$. Simulating that null over 20 trials of 200 memories $\times$
8 tasks gives flip rates of 53.5\%--92.0\% across five settings of $\mu$ and
the memory and task standard deviations. A high flip rate is therefore
consistent with memory value being a scalar plus a task-difficulty term.

\paragraph{The variance decomposition.}
We pre-registered a decomposition $\Delta U(m,t,s) = \mu + a_m + b_t +
(ab)_{mt} + \varepsilon_s$ into component variance shares against
$F_{\mathrm{exec}}$, with a method-of-moments fallback fixed in advance. We
recorded the structural precondition before executing: components are
identified only relative to a residual scale. It fails here because the
reader is deterministic within a session, so we quote no share
(Section~\ref{sec:negative}). Standard errors are clustered by task
throughout.

\paragraph{Similarity and cost.}
Semantic distance is the cosine of the pinned embeddings of memory and task,
banded by terciles fixed before outcomes were examined, since cosine scales
are not comparable across corpora. The seven cost families are
\textbf{storage} (memory tokens held), \textbf{retrieval} (latency, tokens
returned), \textbf{exposure} (tokens placed in the reader's context),
\textbf{generation} (latency, input, output and reasoning tokens),
\textbf{extraction} (formation cost), \textbf{API and compute}, and
\textbf{selectivity} (how much of what was retrieved was worth exposing).
Exposure is recorded separately from retrieval because the gap between them
is truncation loss.
